\twocolumn[
\begin{@twocolumnfalse}
\begin{center}
{\Huge\bfseries A VQE Study of the Lipkin Interaction \par}
\vspace{0.2cm}
{\large FYS5419 -- Quantum computing and quantum machine learning \par}
\vspace{0.2cm}
{\large \textbf{PROJECT 1} \par}
\vspace{0.2cm}
\begin{minipage}{0.48\linewidth}
\raggedright
{\large Egil Furnes \par}
\texttt{egilsf@uio.no} \par
\href{https://github.com/egil10/fys5419}{\texttt{github.com/egil10/fys5419}}
\end{minipage}
\hfill
\begin{minipage}{0.48\linewidth}
\raggedleft
University of Oslo \par
Department of Physics \par
\today
\end{minipage}
\end{center}
\vspace{-1em}
\begin{center}
{\large\bfseries Abstract}
\end{center}

In this project, I study the Lipkin--Meshkov--Glick model using the Variational Quantum Eigensolver (VQE) algorithm on a sequence of increasingly complex quantum systems. Starting from single-qubit operations and Bell state entanglement, I progress through a two-level Hamiltonian with an avoided crossing, a two-qubit interacting system exhibiting interaction-driven entanglement, and finally the $J=1$ and $J=2$ Lipkin Hamiltonians mapped onto two and four qubits respectively. Classical diagonalization using \texttt{numpy.linalg.eigh} achieves machine-precision accuracy ($\approx10^{-15}$) and serves as the reference throughout. Hardware-efficient VQE circuits based on $R_y$ rotations and entangling gates reproduce the exact ground state energies to $\approx10^{-11}$ for the simpler systems. For the $J=2$ Lipkin model the residual VQE error is $\approx10^{-2}$, reflecting the difficulty of optimizing a 16-parameter circuit in a high-dimensional landscape. Hartree-Fock results are included for comparison, and the onset of the quantum phase transition in the Lipkin model is clearly captured in both the classical and VQE spectra.

\vspace{2em}

\end{@twocolumnfalse}
]
